% !TeX root = ../main.tex
\section{Modular Pipelines}
\label{sec:modular}

\subsection{Source separation}
\label{sec:separation}
Transcription models trained on solo guitar degrade sharply on mixes, so a
modular pipeline usually starts by extracting the guitar. Hybrid Transformer
Demucs (HT Demucs)~\cite{rouard2023htdemucs} combines temporal and spectral
U-Nets with a cross-domain transformer and reaches 9.20\,dB SDR on MUSDB18
with extra training data; its six-source variant \texttt{htdemucs\_6s} adds
guitar and piano stems and is the most widely available open guitar
separator. Band-split models operate on subbands of the complex spectrogram:
BS-RoFormer~\cite{lu2023bsroformer} won the music track of the Sound Demixing
Challenge 2023 and reports 9.80\,dB SDR on MUSDB18HQ without extra data, and
Mel-RoFormer~\cite{wang2023melroformer} replaces the heuristic bands with
overlapping mel bands. The MSST framework~\cite{solovyev2026msst} unifies
training and inference for these families, and community checkpoints with a
guitar stem are distributed through it. Query-based separation
(Banquet~\cite{watcharasupat2024banquet}) uses one decoder conditioned on an
instrument query; with 24.9\,M parameters it outperformed
\texttt{htdemucs\_6s} on guitar and piano on MoisesDB and can target narrow
classes such as clean acoustic guitar.

Two limitations matter for tablature. First, these models output one stem for
all guitars: rhythm and lead parts, or two doubled rhythm guitars, come out
together. Spatial cues can partially undo this, since rock guitars are often
panned apart; azimuth discrimination (ADRess)~\cite{barry2004adress} assigns
time-frequency bins to positions in the stereo image. Second, separation
artifacts (bleed from keyboards and synthesizers, smeared transients) are
absent from the solo-guitar data on which transcription models are trained.
A different separation target is the single string: Bastas
et~al.~\cite{bastas2025stringsep} separate a guitar signal into per-string
signals and use them to improve tablature transcription. Tone removal is
another preprocessing option: EG-VAE~\cite{yeh2026egvae} recovers a DI-like
signal from an amplified recording, which could make transcribers trained on
clean guitar applicable to distorted tones.

\subsection{Note transcription}
\label{sec:amt}
\emph{Frame-and-onset models.} Onsets and Frames~\cite{hawthorne2018onsets}
established joint prediction of onsets and frame activity for piano.
Kong et~al.~\cite{kong2021highres} replaced binary onset targets with
regression of the exact onset and offset times, reaching 96.72\% onset F1 on
MAESTRO and robustness to misaligned labels. This architecture is the basis
of current guitar models: Riley et~al.\ adapted it to guitar by aligning
commercially available scores to recordings~\cite{riley2024domain}, obtaining
state-of-the-art zero-shot results on GuitarSet, and then trained it on the
GAPS classical guitar dataset~\cite{riley2024gaps}, again state of the art on
GuitarSet both supervised and zero-shot. Basic Pitch~\cite{bittner2022basicpitch}
is a small instrument-agnostic model that predicts onsets, notes and a pitch
contour at a third of a semitone; it is a common lightweight baseline and its
contour is useful for technique detection, but as an instrument-agnostic
model it tends to report string harmonics as notes.

\emph{Electric guitar.} Timbre is the main obstacle. EGDB showed that
amplifier tone strongly affects accuracy~\cite{chen2022egdb}, and proposed a
multi-loss transformer. EGDB-PG~\cite{chen2025egdbpg} trained a Tone-informed
Transformer (TIT) with a tone embedding on many amplifier-cabinet renders and
reported gains across amplifier types driven by data diversity and the
embedding. Pedroza et~al.~\cite{pedroza2024tones} improved robustness by
training on synthetic performances played through recordings of real guitar
tones and effects. The common recipe is to present the same performance in
many tones, which DI datasets (EGDB, Guitar-TECHS, GOAT) make possible.

\emph{Sequence-to-sequence transcription.} Hawthorne
et~al.~\cite{hawthorne2021seq2seq} showed that a generic encoder-decoder
transformer translating spectrograms into MIDI-like tokens matches
specialised piano models. MT3~\cite{gardner2022mt3} extended this to many
instruments with a shared token vocabulary that includes the instrument, and
reported large gains for low-resource instruments such as guitar.
YourMT3+~\cite{chang2024yourmt3} replaced the encoder with a hierarchical
time-frequency transformer with mixture-of-experts, added a multi-channel
decoder for partially annotated data and cross-dataset stem mixing, and
remained limited on pop recordings. MuScriptor~\cite{rouard2026muscriptor}
departs from architecture changes and scales data instead: a decoder-only
transformer (60\,M to 1.3\,B parameters) reading 5\,s mel-spectrograms is
pre-trained on 1.45\,M synthesised MIDI files, trained on 170,000 aligned real
recordings, and post-trained with a GRPO-style reinforcement learning stage on
300 verified pieces. Distorted electric guitar is the second most frequent
instrument class in its real training set (about 47\% of tracks), and clean
electric and acoustic guitar are separate classes. It can be conditioned on
the instruments present. Its tokenizer cannot represent two simultaneous notes
of the same pitch on the same instrument class; keeping such notes in its test
set lowers onset F1 of the largest model from 60.4 to 51.8. Its weights are
released under a non-commercial license.

For tablature, these models have two shortcomings: they output notes with an
instrument class but no string or fret, and all guitars of the same class are
merged into one stream.

\subsection{Beats, bars and meter}
\label{sec:rhythm}
Writing rhythm requires beats, downbeats and a time signature. Classic
trackers combine an onset function with dynamic programming~\cite{ellis2007beat};
neural trackers predict beat and downbeat activations and decode them with a
dynamic Bayesian network (DBN) that assumes a fixed meter and tempo
range~\cite{bock2016joint}. All-In-One~\cite{kim2023allinone} tracks beats,
downbeats and song sections jointly from separated stems.
Beat This!~\cite{foscarin2024beatthis} removes the DBN, is trained on diverse
data including solo instruments and pieces with time-signature changes, and
surpasses prior work in F1, though it is weaker on continuity metrics.
AlignBeat~\cite{ahn2025alignbeat} models the beat class as latent to train on
partially labelled data and reports 3.6 points higher downbeat CMLt than Beat
This!. Explore This!~\cite{kihlborg2026explore} predicts a tatum grid and
groups it into beats and bars with differentiable bar-pointer modules, which
improves continuity over DBN-free systems. A failure analysis on the SMC
dataset~\cite{ahn2026smc} finds that state-of-the-art trackers produce
``confident but wrong'' activations on expressive music, and that the usual
DBN minimum tempo of 55\,BPM forces double-tempo output on 21\% of SMC tracks.
Quantising notes onto the beat grid and choosing notated durations is rarely
treated in GTT papers; TART uses a rule-based tablature generator for
it~\cite{gupta2026tart}.

\subsection{String and fret assignment}
\label{sec:fingering}
Given notes, a tablature requires choosing among the string-fret positions of
each pitch. \emph{Search-based methods} minimise a cost over the whole
sequence: Sayegh's optimum path paradigm~\cite{sayegh1989fingering} and
hidden Markov models whose states are left-hand forms, decoded with Viterbi
variants~\cite{hori2016minimax}. Their costs (stretch, hand movement,
position) are hand-designed and capture little of the long-range consistency
of human tablature.

\emph{Learned symbolic models} treat fingering as translation from notes to
tablature tokens, pre-trained on DadaGP. MIDI-to-Tab~\cite{edwards2024miditotab}
trains an encoder-decoder transformer with masked language modelling over
string tokens, fine-tunes it on professionally transcribed performances, and
was rated more playable than competing algorithms in a user study.
Fretting-Transformer~\cite{hamberger2025fretting} uses a T5 model trained on
DadaGP, GuitarToday and Leduc, adds tuning and capo conditioning, and
surpasses an A* baseline and Guitar Pro's own conversion. Hsu
et~al.~\cite{hsu2026noteevent} add explicit note-event tokens and a decoding
constraint that masks string-fret candidates inconsistent with the pitch; they
improve on Fretting-Transformer on DadaGP and especially when training only on
the small Leduc set.

\emph{Audio-conditioned fingering.} Symbolic models ignore timbre, which
differs between strings for the same pitch. TART's third stage is a T5
encoder-decoder conditioned on audio features~\cite{gupta2026tart}; together
with its note stage it reaches 71.8\% string-fret Tab F1 averaged over its
benchmarks, 8.5 points above the best prior baseline.

\subsection{Playing techniques}
\label{sec:techniques}
Rule-based detection on the fundamental-frequency contour was introduced by
Chen et~al.~\cite{chen2015technique} for bends, vibrato, slides and
hammer-ons in real electric guitar recordings. TENT~\cite{su2019tent} embeds
technique recognition in note tracking for monophonic solos and reports a
14.7\% higher note F-score than an existing method. TART~\cite{gupta2026tart}
classifies techniques per note with an MLP. GOAT~\cite{loth2025goat} and
IDMT-SMT-Guitar~\cite{kehling2014idmt} provide technique labels for training
and testing. Detection is mostly evaluated on clean, monophonic material;
distortion blurs the pitch contour and compressed attacks hide the pick
transient that separates a picked note from a hammer-on, so techniques on
amplified rhythm parts are largely untested.
